\section{Method}\label{sec:method}

\method{} reconstructs a skeleton-drivable 3D Gaussian avatar from a
monocular video stream in two stages: a \emph{warm-start} that produces a
usable drivable avatar within seconds of the first frame
(Sec.~\ref{sec:method:warmstart}), and a \emph{continual refinement} that
ingests subsequent clips online while mitigating catastrophic forgetting
(Sec.~\ref{sec:method:continual}). Fig.~\ref{fig:pipeline} gives an
overview.

% ----------------------------------
\begin{figure*}[t]
  \centering
  % \includegraphics[width=0.95\textwidth]{figures/pipeline.pdf}
  \fbox{\parbox{0.93\textwidth}{\centering\vspace{1.6em}
    \textbf{[Pipeline placeholder.]} Left: a feed-forward Gaussian
    generator warm-starts a drivable avatar from the first frame together
    with an SMPL registration. Right: the monocular video stream is
    partitioned into clips; each clip triggers a continual-learning step
    combining an EWC prior, a viewpoint-stratified replay buffer, and
    pose-conditioned selective reactivation.
  \vspace{1.6em}}}
  \caption{Overview of \method{}. A drivable 3DGS avatar is warm-started
  from the first frame and then continuously refined as the monocular
  video stream arrives, without forgetting previously learned viewpoints.}
  \label{fig:pipeline}
\end{figure*}

\subsection{Preliminaries}\label{sec:method:prelim}
A 3D Gaussian avatar is represented as a set of $N$ Gaussians, each with
mean $\bm{\mu}\!\in\!\mathbb{R}^3$, scale $\bm{s}\!\in\!\mathbb{R}^3$ (stored
as $\log\bm{s}$), rotation quaternion $\bm{q}\!\in\!\mathbb{R}^4$, opacity
$o$ (stored as logit), and degree-0 spherical-harmonic color $\bm{c}$.
Following articulated 3DGS
avatars~\cite{hu2024gauhuman,hu2024gaussianavatar}, Gaussians are defined
in a canonical (rest-pose) space and deformed to posed space via Linear
Blend Skinning (LBS):
\begin{equation}
  \bm{G} = \sum_{k=1}^{K} w_k\, \bm{G}_k,\quad
  \bm{\mu}^{p} = \bm{G}_{[:3,:3]}\bm{\mu}^{c} + \bm{G}_{[:3,3]},
  \label{eq:lbs}
\end{equation}
where $\{w_k\}$ are per-Gaussian skinning weights to $K$ SMPL
joints~\cite{loper2015smpl} and $\{\bm{G}_k\}$ are per-joint
homogeneous transforms. Rotation and scale are obtained from the polar
decomposition of $\bm{G}_{[:3,:3]}$. Rendering is performed with
differentiable 3DGS rasterization~\cite{kerbl3Dgaussians,ye2024gsplat}.

\subsection{Warm-Start: Fast Drivable Avatar from a Single Frame}\label{sec:method:warmstart}
The warm-start yields a drivable avatar from the first frame
$\mathcal{I}_0$ in seconds, providing a strong initialization for
continual refinement.

\textbf{Feed-forward Gaussian generation.} We apply a pre-trained
feed-forward 3DGS generator $\mathcal{G}$ (TripoSplat~\cite{yan2026triposplat})
to $\mathcal{I}_0$ to obtain an initial set of Gaussians
$\{\bm{\mu}^{c}, \bm{s}, \bm{q}, o, \bm{c}\}$. Background is removed so the
Gaussians concentrate on the subject.

\textbf{Monocular SMPL registration.} A 2D pose estimator
(YOLO11-pose~\cite{ultralytics2023yolo}) extracts keypoints from
$\mathcal{I}_0$, and an SMPL regressor (SMPLer-X~\cite{cai2024smplerx})
produces body shape $\bm{\beta}$, pose $\bm{\theta}_0$, and translation
$\bm{\tau}_0$. Because a single view is ambiguous, we refine the
registration by minimizing a combination of 2D keypoint reprojection and
silhouette loss over a few Adam~\cite{kingma2014adam} steps:
\begin{equation}
  \mathcal{L}_{\mathrm{reg}} = \lambda_{\mathrm{kps}}\mathcal{L}_{\mathrm{kps}}
    + \lambda_{\mathrm{sil}}\mathcal{L}_{\mathrm{sil}}.
  \label{eq:reg}
\end{equation}

\textbf{Skeleton rigging.} Each Gaussian is bound to the SMPL skeleton by
finding its $n$ nearest rest-pose mesh vertices and inheriting their
skinning weights through a distance-softmax:
\begin{equation}
  w_k = \frac{\sum_{j\in\mathcal{N}(\bm{\mu}^{c})}\!\!\exp(-d_j/\tau)\,W_{jk}}
             {\sum_{k'}\sum_{j\in\mathcal{N}(\bm{\mu}^{c})}\!\!\exp(-d_j/\tau)\,W_{jk'}},
  \label{eq:rig}
\end{equation}
where $d_j$ is the distance to vertex $j$, $\tau$ a temperature, and
$W_{jk}$ the SMPL skinning weights. The result is a drivable
avatar $\mathcal{A}_0$ that can be posed by any SMPL parameters---the
minimal usable asset is available immediately.

\subsection{Continual Refinement over Video Streams}\label{sec:method:continual}
We partition the remaining monocular stream into clips
$\{\mathcal{C}_1, \mathcal{C}_2, \dots\}$ of $T$ frames each. Each clip is
treated as a continual-learning \emph{task}: we optimize the avatar on
$\mathcal{C}_t$ while preserving knowledge from
$\{\mathcal{C}_1,\dots,\mathcal{C}_{t-1}\}$. For each clip we estimate
SMPL parameters per frame (as in the warm-start), and optimize the
canonical Gaussian parameters $\bm{\Theta}=\{\bm{\mu}^{c},\log\bm{s},\bm{q},\mathrm{logit}\,o,\bm{c}\}$
(and optionally the skinning-weight logits) by deforming to each frame's
pose and rendering against the observed image.

\subsubsection{Base reconstruction loss}
For a frame $\mathcal{I}$ with estimated pose $\bm{\theta}$ and camera
$\bm{\pi}$, the rendered image $\hat{\mathcal{I}}$ is supervised with the
standard 3DGS photometric loss:
\begin{equation}
  \mathcal{L}_{\mathrm{rec}} = \lambda_1 \|\hat{\mathcal{I}}-\mathcal{I}\|_1
    + \lambda_{\mathrm{ssim}}\bigl(1-\mathrm{SSIM}(\hat{\mathcal{I}},\mathcal{I})\bigr).
  \label{eq:rec}
\end{equation}

\subsubsection{EWC prior for Gaussian parameters}\label{sec:method:ewc}
Naively minimizing $\mathcal{L}_{\mathrm{rec}}$ on $\mathcal{C}_t$ causes
catastrophic forgetting of Gaussians well-constrained by earlier clips.
Inspired by Elastic Weight Consolidation~\cite{kirkpatrick2017overcoming},
we estimate, after finishing clip $\mathcal{C}_{t-1}$, a per-parameter
importance $F_i$ as the empirical Fisher information, \ie{} the squared
gradient of $\mathcal{L}_{\mathrm{rec}}$ accumulated over the clip's frames
and over a small held-out replay subset:
\begin{equation}
  F_i \leftarrow \gamma\, F_i + (1-\gamma)\,
  \frac{1}{|\mathcal{C}_{t-1}|}\sum_{\mathcal{I}\in\mathcal{C}_{t-1}}
  \Bigl(\partial_{\Theta_i}\mathcal{L}_{\mathrm{rec}}\Bigr)^{2},
  \label{eq:fisher}
\end{equation}
where $\gamma\!\in\![0,1]$ is a forgetting factor for the running
importance and $\Theta_i$ denotes a single scalar parameter. The
importance $F_i$ is large for parameters that strongly affect the
rendering of previously observed viewpoints---typically Gaussians that are
visible and well-constrained in those views. When learning from the new
clip $\mathcal{C}_t$, we add a quadratic penalty anchoring parameters to
their post-$\mathcal{C}_{t-1}$ values $\Theta^{\star}_i$:
\begin{equation}
  \mathcal{L}_{\mathrm{ewc}} = \sum_{i} \tfrac{1}{2}\, F_i\,(\Theta_i - \Theta^{\star}_i)^{2}.
  \label{eq:ewc}
\end{equation}
Because $F_i$ is defined per Gaussian attribute, the penalty acts
spatially: it protects well-observed body regions (high $F_i$) while
leaving under-observed ones (low $F_i$) free to update, naturally
allocating the model's plasticity to where new evidence arrives.

\subsubsection{Viewpoint-stratified replay buffer}\label{sec:method:replay}
EWC alone is a soft constraint and can still drift over many clips. We
complement it with an episodic-memory replay buffer $\mathcal{B}$ that
stores a compact, pose-diverse set of historical frames. To keep $\mathcal{B}$
representative of the articulated body's viewpoint distribution, we
stratify by SMPL global orientation: frames are binned by the azimuth of
the root pose, and within each bin a fixed quota is retained using a
coverage-greedy selection that maximizes the minimum feature distance to
already-buffered frames (herding~\cite{rebuffi2017icarl}). At each
optimization step on $\mathcal{C}_t$, we additionally render and supervise
a small batch sampled from $\mathcal{B}$, yielding a replay loss
$\mathcal{L}_{\mathrm{rb}}$ with the same form as
Eq.~\eqref{eq:rec}. This directly re-exposes the optimizer to earlier
viewpoints and provides a stronger anti-forgetting signal than the
parameter-space EWC prior.

\subsubsection{Pose-conditioned selective reactivation}\label{sec:method:react}
Streaming clips are pose-biased: a clip may show only the upper body, the
back, or a particular garment region. To use the optimization budget
efficiently and avoid densifying already-well-covered regions, we employ
a selective reactivation policy. For each Gaussian we maintain an
accumulated visibility count $v_i$ over observed frames. During the
densification phase of 3DGS~\cite{kerbl3Dgaussians}, candidates for
clone/split are filtered so that Gaussians in body regions whose SMPL
joints have low historical visibility are prioritized; regions with high
$v_i$ are pruned more aggressively. Concretely, we modulate the
densification gradient threshold per Gaussian by
\begin{equation}
  g_i = g_0 \cdot \bigl(1 + \rho\, \bar{v}_i\bigr),
  \label{eq:react}
\end{equation}
where $\bar{v}_i$ is the normalized visibility and $\rho$ controls the
strength. This raises the split/clone bar for frequently seen Gaussians
and lowers it for under-observed ones, steering capacity toward novel
detail as the stream progresses.

\subsubsection{Overall objective}
The total loss during clip $\mathcal{C}_t$ is
\begin{equation}
  \mathcal{L} = \mathcal{L}_{\mathrm{rec}}(\mathcal{C}_t)
              + \lambda_{\mathrm{ewc}}\mathcal{L}_{\mathrm{ewc}}
              + \lambda_{\mathrm{rb}}\mathcal{L}_{\mathrm{rb}}(\mathcal{B})
              + \lambda_{\mathrm{w}}\mathcal{L}_{\mathrm{w}},
  \label{eq:total}
\end{equation}
where $\mathcal{L}_{\mathrm{w}}$ is an optional L2 regularizer keeping
optimizable skinning weights close to the SMPL prior~\cite{hu2024gauhuman}.
After $\mathcal{C}_t$ is processed, the Fisher map $F_i$ (Eq.~\eqref{eq:fisher}),
the replay buffer $\mathcal{B}$, and the visibility counts $v_i$ are
updated, and the avatar is immediately usable for rendering and animation.
The procedure is summarized in Algorithm~\ref{alg:continualgs}.

\begin{algorithm}[t]
\caption{\method{} continual refinement}
\label{alg:continualgs}
\begin{algorithmic}[1]
\Require Warm-start avatar $\mathcal{A}_0$, stream clips $\{\mathcal{C}_t\}$
\State Initialize Fisher map $F\!\leftarrow\!\bm{0}$, replay buffer $\mathcal{B}\!\leftarrow\!\emptyset$, visibility $v\!\leftarrow\!\bm{0}$
\For{each clip $\mathcal{C}_t$ in the stream}
  \State Estimate per-frame SMPL pose $\bm{\theta}$ and camera $\bm{\pi}$
  \For{a fixed number of optimization steps}
    \State Sample batch from $\mathcal{C}_t$ and from $\mathcal{B}$
    \State Deform canonical Gaussians via LBS (Eq.~\eqref{eq:lbs}); render
    \State Compute $\mathcal{L}$ (Eq.~\eqref{eq:total}) and backpropagate
    \State Densify with pose-conditioned thresholds (Eq.~\eqref{eq:react})
  \EndFor
  \State Update $F$ via accumulated squared gradients (Eq.~\eqref{eq:fisher})
  \State Update $\mathcal{B}$ via viewpoint-stratified herding
  \State Update visibility counts $v$; emit drivable avatar $\mathcal{A}_t$
\EndFor
\end{algorithmic}
\end{algorithm}

\subsection{Discussion}\label{sec:method:discussion}
\method{} decouples \emph{availability} from \emph{fidelity}: the
warm-start avatar is drivable immediately, and each subsequent clip
improves a specific aspect of the body under the protection of the EWC
prior and replay buffer. The design is orthogonal to the choice of
feed-forward generator and SMPL regressor, so improvements in either
component directly benefit \method{}. Compared with streaming 3DGS
methods~\cite{wu2025streamsplat,li2025streamgs} that do not model
articulated structure, our skeleton-conditioned deformation ensures the
avatar remains drivable throughout; compared with per-video
optimizers~\cite{hu2024gauhuman,hu2024gaussianavatar,qian2024gsavatar},
the continual-learning formulation enables online, incremental consumption
of footage without retraining.
